\chapter{Real Structures on Wallpaper Groups} \label{appendix_wallpaper}

In this appendix, we rederive the classification of magnetic wallpaper groups, using our result on the classification of Real structures from Chapter \ref{chapter:realstructures}. We also give a second proof of the classification, by explicitly counting index two subgroups of wallpaper groups. We use this explicit presentation when we compute the antilinear representation theory of these groups, in Appendix \ref{appendix_repswallpaper}.

\section{Wallpaper Groups}

\subsection{Presentation} Let us first recall the classification of wallpaper groups, that is, $2$-dimensional crystallographic groups, as considered in Chapter \ref{chap:physics}.

We give generators for each wallpaper group inside $\Euc_2 = \HO_2 (\bR)\ltimes \bR^2$. The following matrices are useful. Their index denotes their order.
\[
I =\begin{pmatrix}1&0\\0&1\end{pmatrix},\quad
U_2=\begin{pmatrix}0&1\\1&0\end{pmatrix},\quad
W_2=\begin{pmatrix}-1&0\\0&-1\end{pmatrix},\quad
X_2=\begin{pmatrix}1&0\\0&-1\end{pmatrix},\quad
X_3=\begin{pmatrix}-1&-1\\1&0\end{pmatrix},
\]

\[
X_4=\begin{pmatrix}0&-1\\1&0\end{pmatrix}, \quad
X_6=\begin{pmatrix}0&-1\\1&1\end{pmatrix},\quad
Y_2=\begin{pmatrix}-1&0\\0&1\end{pmatrix}, \quad
Z_2=\begin{pmatrix}0&-1\\-1&0\end{pmatrix}, \quad
X_{\infty}=\begin{pmatrix}1&1\\0&1\end{pmatrix}.
\]

Recall that $[A\mid x]$ denotes the affine transformation $y\mapsto Ay+x$, for $A \in \GL_2(\bR)$ and $x,y \in \bR^2$, column vectors. Any element $v \in \bR^2$ defines a translation $[I_2 \mid v]$. In particular, we will implicitly $\bR^2$ as a subset of $\Euc_2$. Let $e_1$ and $e_2$ denote the standard basis of $\bR^2$. In Table \ref{table:wallpaper1}, we list the 17 isomorphism classes of wallpaper groups, each such $G$ given by generators, so that the lattice it defines $T \coloneqq G \cap \bR^2$ is the subgroup $\langle e_1, e_2 \rangle$. 

\begin{table}[!htbp]
  \centering
\small
\begin{tabular}{lllllll}
\toprule
\textbf{Group} & \textbf{Orbifold} & \textbf{Cell} & \textbf{Centre} & $\boldsymbol{Q}$ & $\boldsymbol{G}=\langle e_1,e_2,$ & $\boldsymbol{Z(G)}$ \\
\midrule
p1   & $o$       & oblique     &  & $C_1=\langle \emptyset \rangle$ &  $\emptyset \rangle$ & $T=\langle e_1,e_2\rangle$ \\
p2   & $2222$    & oblique     &  & $C_2=\langle W_2\rangle$ &  $W_2\rangle$ & $0$ \\
\midrule
pm   & $**$      & rectangular &  & $C_2=\langle X_2\rangle$ & $X_2\rangle$ & $\bZ =\langle e_1\rangle$ \\
pg   &  $xx$      & rectangular &  & $C_2=\langle X_2\rangle$ & $[X_2\mid \tfrac12 e_1]\rangle$ & $\bZ=\langle e_1\rangle$ \\
cm   & $*x$      & rhombic & $\tfrac12 (e_1+e_2) $ & $C_2=\langle U_2\rangle$ & $U_2\rangle$ &  $\bZ =\langle e_1+e_2\rangle$ \\
pmm  & $*2222$   & rectangular &  & $K_4=\langle X_2,Y_2\rangle$ & $X_2,Y_2\rangle$ & $0$\\
pmg  &  $22*$     & rectangular &  & $K_4=\langle X_2,Y_2\rangle$ &$X_2, [Y_2\mid\tfrac12 e_2] \rangle$ &  $0$\\
pgg  & $22x$     & rectangular &  & $K_4=\langle X_2,Y_2\rangle$ & $[X_2\mid\tfrac12 e_1] , [Y_2\mid\tfrac12 e_2] \rangle$ & $0$ \\
cmm  &  $2\hspace{-0.25em}*\hspace{-0.25em}22$    & rhombic & $\tfrac12 (e_1+e_2) $ & $K_4=\langle U_2,W_2\rangle$ &$U_2,W_2\rangle$ & $0$ \\
\midrule
p4   & $442$     & square      &  & $C_4=\langle X_4\rangle$ &  $X_4\rangle$ &   $0$  \\
p4m  & $*442$    &  square      & & $D_4=\langle X_4,X_2\rangle$ &  $X_4,X_2\rangle$ &  $0$   \\
p4g  & $4\hspace{-0.25em}*\hspace{-0.25em}2$     &  square      & & $D_4=\langle X_4,X_2\rangle$ &  $X_4,[X_2\mid\tfrac12 (e_1+e_2)]\rangle$ &  $0$   \\
\midrule
p3   & $333$     &  hexagonal   &  & $C_3=\langle X_3\rangle$ &  $X_3\rangle$ &  $0$  \\
p3m1 & $*333$    & hexagonal   & & $D_3=\langle X_3,U_2\rangle$ &  $X_3,U_2\rangle$ &   $0$  \\
p31m & $3\hspace{-0.25em}*\hspace{-0.25em}3$     & hexagonal   & & $D_3=\langle X_3,Z_2\rangle$ &  $X_3,Z_2 \rangle$ &   $0$   \\
p6   &  $632$     & hexagonal   &  & $C_6=\langle X_6\rangle$ &$X_6\rangle$ &  $0$  \\
p6m  & $*632$    & hexagonal   &  & $D_6=\langle X_6,U_2\rangle$ &  $X_6,U_2\rangle$ &  $0$   \\
\bottomrule
\end{tabular}
\caption{The 17 wallpaper groups with generators, orbifold notation and centres.}
\label{table:wallpaper1}
\end{table}

Let us explain the rest of the entries of the table.

\begin{itemize}
  \item  The cell is the parallelogram spanned by $e_1,e_2$, for which we also record the centre of the cell. The type of the cell corresponds to restriction on the basis.
Namely, oblique means no restriction,
rectangular  -- $e_1,e_2$ are orthogonal,
rhombic  -- $e_1,e_2$ are of the same length,
square --  $e_1,e_2$ is an orthonormal basis,
hexagonal -- $e_1,e_2$ is a pair of vectors of the same length
with the interior angle having 60 degrees. 
\item $Z(G)$ is the centre of $G$, and $Q$ is the point group $G / T$ of $G$.
\end{itemize}

Note that there are four conventional notations for the wallpaper groups. We are not using the spectroscopy (Schoenflies) and algebra (Coxeter) notations. Our main notation is the crystallography (Hermann-Mauguin) notation. It is essentially a recipe for both constructing the group and recognising it from the pattern:
\begin{itemize}
\item The first letter stands for {\bfseries p}rimitive or {\bfseries c}entred cell. 
  A centred cell contains a centre listed in column 4; a primitive cell does not.
\item $n\in \{1,2, 3, 4,6\}$ stand for a rotation symmetry by $\frac{360}{n}$ degrees.
\item {\bfseries m} stands for mirror (a reflection is required to generate $G$).
\item {\bfseries g} stands for glide (a glide reflection is required to generate $G$).
\item Both {\bfseries p3m1} and {\bfseries p31m} are semidirect products of $D_3=S_3$ with $T$ but the representations
  are not equivalent. Interestingly enough, if one embeds $S_3$ into $\GL_2 (\bZ)$, one has two representations on $\bZ^2$, one
  on rows, another on columns. The corresponding semidirect products will be {\bfseries p3m1} and {\bfseries p31m}.
  For another algebraic difference, observe
  $H^1(Q,T)$ in Table~\ref{table:wallpaper2}.
\item There are also two geometric differences between {\bfseries p3m1} and {\bfseries p31m}
that help to distinguish them visually.
One difference is that  in {\bfseries p3m1} the mirrors lie between the vectors of $T$, while
in {\bfseries p31m} the mirrors run parallel to the vectors in $T$.
The second difference is that
 in {\bfseries p3m1} all the centres of 3-fold reflections lie on the mirrors, while
 in {\bfseries p31m} there are two types of 3-fold rotation centres: lying on the mirrors and not lying on the mirrors.
\end{itemize} 

We also use the orbifold (Conway) notation. The notation is simply a list of all singular points on the orbifold $\bR^2/G$.
\begin{itemize}
\item $o$ stands for the empty set $\emptyset$, i.e., there are no singular points.
\item 2, 3, 4 and 6 stand for a cone point of the corresponding order.
\item $*$ stands for a boundary component: a connected set of mirror points.
\item 2, 3, 4 and 6 after $*$ means that the cone point lies on the boundary.  
\item $x$ stands for a boundary component: a connected set of glide points.
\end{itemize}   

\subsection{Outer Automorphism Groups} Let us now consider the outer automorphism groups of the 17 wallpaper groups, these being needed to compute their Real structures via the methods of Chapter \ref{chapter:realstructures}. We record these, and the associated data we need to compute them in Table \ref{table:wallpaper2}. We give explicit generators for these outer automorphism groups in Table \ref{table:wallpaper2_geb}.

\begin{table}[!htbp]
\centering
\small
\begin{tabular}{lllllllll}
\toprule
 $\boldsymbol{G}$ &  $\boldsymbol{Q}$ &  $\boldsymbol{H^1(Q,T)}$  & $\boldsymbol{\Out(Q)}$ & $\boldsymbol{C_{\GL_2 (\bZ) }(Q)}$ & $\boldsymbol{M}$  & $\boldsymbol{N(Q)}$ & $\boldsymbol{N(Q,o)}$ &  $\boldsymbol{\Out(G)}$  \\
\midrule
p1   & $C_1$       & $0$      &  $C_1$ & $\GL_2 (\bZ)$  &   & $\GL_2 (\bZ)$ &  $\GL_2 (\bZ)$ &  $\GL_2 (\bZ)$ \\
p2   & $C_2$    & $(\bZ/2)^2$  & $C_1$  &  $\GL_2 (\bZ)$  &   &  $\PGL_2 (\bZ)$ & $\PGL_2 (\bZ)$ & $\PGL_2 (\bZ)\ltimes K_4$  \\
\midrule
pm   & $C_2$      &  $\bZ/2$ & $C_1$  & $K^a_4$ &   & $C_2$ & $C_2$ &  $K_4$ \\
pg   &  $C_2$      & $\bZ/2$  & $C_1$  & $K^a_4$ &   & $C_2$ & $C_2$ &  $K_4$\\
cm   & $C_2$      & $0$  & $C_1$  & $K^b_4$ &   & $C_2$ & $C_2$ &  $C_2$ \\
pmm  & $K_4$   & $(\bZ/2)^2$  & $\boldsymbol{D_3}$  & $K_4^a$ & $X_4$ & $C_2$ &$C_2$ &  $D_4$\\
pmg  &  $K_4$     & $(\bZ/2)^2$  & $\boldsymbol{D_3}$  & $K_4^a$ & $X_4$ &   $C_2$ & $C_1$ &  $K_4$\\
pgg  & $K_4$     & $(\bZ/2)^2$  & $\boldsymbol{D_3}$  & $K_4^a$ & $X_4$ &    $C_2$ & $C_2$ &  $D_4$\\
cmm  &  $K_4$    & $\bZ/2$ & $\boldsymbol{D_3}$  &  $K^b_4$ & $X_4$ &    $C_2$ & $C_2$ &  $K_4$\\
\midrule
p4   & $C_4$     & $\bZ/2$       & $C_2$ & $C^a_4$  & $X_2$ &  $C_2$  &  $C_2$ &  $K_4$\\
p4m  & $D_4$    &  $\bZ/2$       & $\boldsymbol{C_2}$ & $C^a_2$ &   &  $C_1$ &  $C_1$ &  $C_2$   \\
p4g  & $D_4$     & $\bZ/2$      & $\boldsymbol{C_2}$ & $C^a_2$  &   &  $C_1$ &  $C_1$ &  $C_2$   \\
\midrule
p3   & $C_3$     & $\bZ/3$    & $C_2$ &  $C^a_6$ & $U_2$ &  $K_4$ & $K_4$ &  $D_6$\\
p3m1 & $D_3$    & $\bZ/3$   & $C_1$ & $C^a_2$ &   & $C_2$ &  $C_2$ &  $D_3$   \\
p31m & $D_3$     & $0$    & $C_1$ & $C^a_2$  &   & $C_2$ &  $C_2$ &  $C_2$  \\
p6   &  $C_6$     & $0$       & $C_2$ &  $C^a_6$ & $U_2$       & $C_2$ &  $C_2$ &  $C_2$  \\
p6m  & $D_6$    & $0$   & $\boldsymbol{C_2}$ & $C^a_2$ &   &  $C_1$ &  $C_1$ &  $C_1$   \\
\bottomrule
\end{tabular}
\caption{The outer automorphisms of the 17 wallpaper groups.}
\label{table:wallpaper2}
\end{table}

\begin{table}[!htbp]
\centering
\small
\begin{tabular}{ll|l|l}
\toprule
 $\boldsymbol{G}$ & $\boldsymbol{\Out(G)}$  & $A, \ B, \ C$ &  Relations of $\boldsymbol{\Out (G)}$ \\
\midrule
p1   & $\GL_2 (\bZ)$  & $X_2, \, X_3, \, X_4$ &  $A^2$, $B^3$, $C^4$, $(AC)^2$, $(CAB)^2$, $[C^2,B]$ \\
p2   & $\PGL_2 (\bZ)\ltimes K_4$   & $[I \mid\tfrac12 e_1], \ U_2, \ X_{\infty}$ & $A^2$, $B^2$, $[A,BAB]$,  $[A,C]$, $CBABC^{-1}(AB)^{-2}$, \\
   &   &  & $(C^{-2}BCB)^2$, $(BCBC^{-1})^3$ \\
\midrule
pm   & $K_4$  &  $W_2, \ [I \mid\tfrac12 e_2]$ & $A^2$, $B^2$, $[A,B]$ \\
pg   & $K_4$ &  $Y_2, \ [I \mid\tfrac12 e_2]$ & $A^2$, $B^2$, $[A,B]$ \\
cm   & $C_2$  & $W_2$ & $A^2$ \\
pmm  & $D_4$ &  $U_2, \ [I \mid\tfrac12 e_2]$ & $A^2$, $B^2$, $(AB)^4$ \\
pmg  & $D_4$ &  $[I \mid\tfrac12 e_1], \ [I \mid\tfrac12 e_2]$ & $A^2$, $B^2$, $[A,B]$ \\
pgg  & $D_4$ &  $U_2, \ [I \mid\tfrac12 e_2]$ & $A^2$, $B^2$, $(AB)^4$ \\
cmm  & $K_4$ &  $X_4, \ [I \mid\tfrac12 (e_1+ e_2)]$ &  $A^2$, $B^2$, $[A,B]$ \\
\midrule
p4   &   $K_4$ & $X_2, \ [I \mid\tfrac12 (e_1+ e_2)]$ &  $A^2$, $B^2$, $[A,B]$ \\
p4m  &  $C_2$    & $[I \mid\tfrac12 (e_1+ e_2)]$ & $A^2$ \\
p4g  &  $C_2$    & $[I \mid\tfrac12 (e_1+ e_2)]$ & $A^2$ \\
\midrule
p3   &  $D_6$   & $W_2, \ U_2, \ [I \mid\tfrac13 (e_1+ e_2)]$ &  $A^2$, $B^2$, $C^3$, $[A,B]$, $[B,C]$, $(AC)^2$ \\
p3m1 &  $D_3$    &  $X_6, \ [I \mid\tfrac13 (e_1+ e_2)]$ &  $A^2$, $B^3$, $(AB)^2$ \\
p31m &  $C_2$   &  $X_6$ &  $A^2$ \\
p6   &  $C_2$   & $U_2$ & $A^2$ \\
p6m  &  $C_1$    &  &  none \\
\bottomrule
\end{tabular}
\caption{Generators of $\Aut (G)$ and relations of $\Out (G)$ for the 17 wallpaper groups.}
\label{table:wallpaper2_geb}
\end{table}

As each point group $Q$ is finite, we may immediately write down $\Aut(Q)$ and $\Out (Q)$ for the points groups $Q$ for all the wallpaper groups.
To compute the normaliser $N_{\GL_2 (\bZ) }(Q)$, we list the centraliser $C_{\GL_2 (\bZ) }(Q)$ and an additional matrix $M$ so that
\[ N_{\GL_2 (\bZ) }(Q) = \langle Q, C_{\GL_2 (\bZ) }(Q), M \rangle \, .\]
Note that  $N_{\GL_2 (\bZ) }(Q) / C_{\GL_2 (\bZ) }(Q)$ is a subgroup of $\Out (Q)$:
in Table~\ref{table:wallpaper2} we make $\Out (Q)$ bold in those cases when it is a proper subgroup.
There we use the following subgroups of $\GL_2 (\bZ)$:
\[
K^a_4 = \langle X_2,Y_2\rangle = \begin{pmatrix} \ast & 0 \\ 0 & \ast \\ \end{pmatrix}, \
K^b_4 = \langle U_2,W_2\rangle = \begin{pmatrix} x & y \\ y & x \\ \end{pmatrix}, \
C^a_2=\langle W_2\rangle = \begin{pmatrix} x & 0 \\ 0 & x \\ \end{pmatrix},
\]
\[
C^a_4=\langle X_4\rangle = \begin{pmatrix} x & y \\ -y & x \\ \end{pmatrix}, \
C^a_6 = \langle X_6\rangle = \begin{pmatrix} x & y \\ -y & x+y \\ \end{pmatrix}.
\]
The short exact sequence~\eqref{Physics_OutSequence} always splits in the dimension 2,
 so that $\Out (G)$ is computed for Table~\ref{table:wallpaper2} as
\begin{equation} \label{table:OutG}
\Out (G) = N(Q,o) \ltimes H^1(Q,T) \, .
\end{equation}
We also tabulate the companion group
$N (Q) \coloneqq N_{\GL_2 (\bZ) }(Q)/Q$,
which turns out to be equal to $N(Q,o)$ in all cases except {\bfseries pmg}.

An inquisitive reader may have further questions about~\eqref{table:OutG}.
We know it is split automatically, only if the sequence $1 \to T \to G \to Q \to 1$ is split. 
This is true in the thirteen cases except
{\bfseries pg},
{\bfseries pmg},
{\bfseries pgg}
and
{\bfseries p4g}.
In two of them, 
{\bfseries pg}
and
{\bfseries p4g},
the sequence~\eqref{table:OutG} collapses because one of the groups is trivial
but the other two still present a challenge.
Moreover, in five of the groups, namely,
{\bfseries pmm},
{\bfseries pmg},
{\bfseries pgg},
{\bfseries p3}
and
{\bfseries p31m},
there are two possible actions of
$N_{\GL_2 (\bZ) }(Q)/Q)$ on  $H^1(Q,T)$, trivial and non-trivial.
In all five case, the action is non-trivial.

One method to determine the action of $N_{\GL_2 (\bZ) }(Q)/Q$ on $H^1(Q,T)$ and the structure of $\Out (G)$ is to use~\eqref{formula:AUT1},
which is essentially the second Bieberbach Theorem. We carry this out in Table~\ref{table:wallpaper2_geb}. In the third column we list
$A,B,C \in \Aff_2$ so that
\[ \Aut (G) = \langle \gamma_A,\gamma_B,\gamma_C,\gamma_g \,\mid\, g\in G \rangle \quad \mbox{ where } \
\gamma_x (y) = xyx^{-1} \, .\]
If we need less then three generators, we use only $A$ or $A,B$. The group $G$ is always generated as in Table~\ref{table:wallpaper1}.
The last column are relations in $\Out (G)$. For instance,
the first row tells us that in the type {\bfseries p1} we have isomorphisms
\[
\Out(G) \cong \GL_2 (\bZ) \cong \langle A,B,C \, \mid \, A^2 = B^3 = C^4 = (AC)^2= (CAB)^2 = C^2BC^{-2}B^{-1} = 1 \rangle \,
\]
and the quotient map $\Aut (G) \to \Out (G)$ is given by
\[ \gamma_{X_2} \mapsto A, \
\gamma_{X_3} \mapsto B, \
\gamma_{X_4} \mapsto C. \] 

\section{Index Two Subgroups of Wallpaper Groups}

In Table~\ref{table:wallpaper3} we list the index two subgroups of each wallpaper group, grouped by $\Aut (G)$-orbit.
For instance, 3p1 in the first row means 3 subgroups of type p1, which form a single $\Aut(G)$-orbit.
Of course, we simply write p1 instead of 1p1. This is a complete list of the black-and-white magnetic groups.

A subgroup $H$ of index 2 is necessarily necessarily normal and $H\cap T$ has index at most 2 in $T$.
By {\em the translation index} or simply t-index, we understand this index $|T:T\cap H|$, and we group subgroups by this index into type 1 and type 2 subgroups. Note that any subgroup $H$ of index 2 contains $2T$. Therefore, the subgroups of index 2 in $G$ and $G/2T$ are in one-to-one correspondence. When identifying the group $G / 2T$, we use $[n,k]$ to denote SmallGroup(n,k).
Notice that this group is finite and fits into an exact sequence $1\to C_2^n \to G/2T \to Q \to 1$.

An easy method to work out the subgroups of index 2 and $G/2T$ is to use our explicit presentation of $G$ found in the magma code at \cite{GITHUB}. This also makes the presentation of $G/2T$ immediate. The index 2 subgroups are kernels of homomorphism to $C_2$, which one lists from the presentations.

For example, in type {\bfseries p4m} the presentation is
\[
\langle s,t,x,y \mid
    s^2,
    t^4,
    t^s = t^3,
    x^s = x,
    y^s = y^{-1},
    x^t = y^{-1},
    y^t = x,
    y^x = y
\rangle
\]
where $x,y$ are the basis of $T$, and for example $x^t = t^{-1} x t$. Homomorphisms $\phi: G \to C_2$ correspond to choices of $a,b,c\in C_2$ so that
\[ \phi (r) = a, \ \phi (s) = b, \ \phi (x) = c, \ \phi (y) = c. \] 
This gives 7 non-trivial homomorphisms and thus 7 subgroups of index 2.

\begin{table}[!htbp]
\centering
\small
\begin{tabular}{llll|llll}
\toprule
 $\boldsymbol{G}$ &  $\boldsymbol{G/2T}$ &  type 1  & type 2 & $\boldsymbol{\sN}$ & $\boldsymbol{\Theta(G)}$ & $\boldsymbol{Z(G)_{\theta}}$ & $\boldsymbol{\fR_{\theta}(G)}$ \\
\midrule
p1   & $K_4 = [4,2]$       &      &  $3\mbox{p1}$ & 6  & $I$  & $\bZ^2_{+}$ & 1+3 \\
&       &      &   &   & $U_2$  & $\bZ_{d}$  & 1 \\
&       &      &   &   & $X_2$  & $\bZ_{+}\oplus \bZ_{-}$  &  1+1 \\
&       &      &   &   & $W_2$  & $\bZ^2_{-}$  & 1 \\
\midrule
p2   & $C_2^3 = [8,5]$    & $\mbox{p1}$  & $6\mbox{p2}$  &  7  &  7 & &  1 \\
\midrule
pm   & $C_2^3 = [8,5]$     &  $\mbox{p1}$ & $\mbox{pm},\mbox{pg},2\mbox{pm},2\mbox{cm}$ & 6  & $I$ & $\bZ_{+}$  & 1+1 \\
&       &      &   &   & $W_2$  & $\bZ_{-}$  & 1 \\
&       &      &   &   & $[I \mid\tfrac12 e_2]$  & $\bZ_{+}$  & 1+1 \\
&       &      &   &   & $[W_2 \mid\tfrac12 e_2]$  & $\bZ_{-}$  & 1 \\
\midrule
pg   & $C_2\times C_4 = [8,2]$     & $\mbox{p1}$  & $2\mbox{pg}$ & 6 &  $I$ & $\bZ_{+}$  & 1+1 \\
&       &      &   &   & $Y_2$  & $\bZ_{-}$  & 1 \\
&       &      &   &   & $[I \mid\tfrac12 e_2]$  & $\bZ_{+}$  & 1+1 \\
&       &      &   &   & $[Y_2 \mid\tfrac12 e_2]$  & $\bZ_{-}$  & 1 \\
\midrule
cm   & $D_8 = [8,3]$    & $\mbox{p1}$  & $\mbox{pm}, \mbox{pg}$  & 3  & $I$ & $\bZ_{+}$  & 1+1 \\
&       &      &   &   & $W_2$  & $\bZ_{-}$  & 1 \\
\midrule
pmm  & $C_2^4 =[16,14]$   & $\mbox{p2},2\mbox{pm}$  & $4\mbox{pmm},4\mbox{pmg},4\mbox{cmm}$      & 4  & 4  & & 1 \\
pmg  & $C_2\times D_4 =[16,11]$   & $\mbox{p2},\mbox{pm},\mbox{pg}$ & $2\mbox{pmg},2\mbox{pgg}$  & 4  & 4  & & 1  \\
pgg  & $C_4\ltimes K_4 =[16,3]$   & $\mbox{p2},2\mbox{pg}$          &                            & 4  & 4  & & 1  \\
cmm  &  $C_2\times D_4 =[16,11]$    & $\mbox{p2},2\mbox{cm}$ & $2\mbox{pmg},\mbox{pmm},\mbox{pgg}$ & 4 & 4 & & 1 \\
\midrule
p4   & $C_4\ltimes K_4 =[16,3]$      & $\mbox{p2}$  & $2 \mbox{p4}$    & 4  & 4 &  &  1  \\
p4m  & $K_4 \wr C_2 =[32,27]$     & $\mbox{p4},\mbox{pmm},\mbox{cmm}$  & $2\mbox{p4m},2\mbox{p4g}$       & 2  & 2 &  & 1  \\
p4g  & $C_4 \ltimes C_2^3 =[32,6]$    & $\mbox{p4},\mbox{cmm},\mbox{pgg}$  &       & 2  & 2 &  &  1 \\
\midrule
p3   & $A_4 = [12,3]$     &    &  &  4  & 4 &  &  1 \\
p3m1 & $S_4= [24,12]$     & $\mbox{p3}$   &  & 2 & 2 &  & 1 \\
p31m & $S_4= [24,12]$    & $\mbox{p3}$    & & 2  & 2 & &  1 \\
p6   & $C_2\times A_4 =[24,13]$     & $\mbox{p3}$       &  &  2 & 2       &  &  1  \\
p6m  & $C_2 \times S_4 = [48,48]$    & $\mbox{p6},\mbox{p3m1},\mbox{p31m}$   &  & 1 & 1  & &  1 \\
\bottomrule
\end{tabular}
\caption{The real structures of the 17 wallpaper groups.}
\label{table:wallpaper3}
\end{table}


\subsection{Real Structures}

In the right half of Table~\ref{table:wallpaper3} we collect all the data related to the classification of Real structures. As explained in Section \ref{sect:are_there}, the equivalence relation we are considering on magnetic groups corresponds to the weak-equivalence on Real structures considered in Section \ref{sect:weakequiv}. From Lemma ~\ref{lemma:weak_eq}, we have that 
\begin{equation} \label{Wallpaper:RTheta}
  \fR(G)/\overset{\text{magn}}{\sim}  = \bigsqcup_{\theta \in \Theta(G)} \Out_{\theta}(G) \backslash \fR_{\theta}(G) \, , \quad
  \Theta (G) \coloneqq \Out(G) \backslash \{ \theta \in \Out(G)\,\mid\, \theta^2 =1\}\, .
\end{equation}

The sets  The sets $\fR_{\theta}(G)$ are also easier in the crystallographic case because $Z(G)=T^Q$ is always torsion-free,
hence, classified by Proposition~\ref{prop:torsionfreeKS}.
Furthermore, $Z(G)=0$ in most  of the cases.

The column ${\sN}$ simply lists the size of the left-hand set in~\eqref{Wallpaper:RTheta}.
It can be counted using Fact~\ref{sect:are_there}.(I).
Indeed, ${\sN}$ is simply equal to 1 plus the number of times the corresponding orbit
appears columns 3 and 4.
For instance, we can see that {\bfseries pmm}
appears once in column 3 and two times in column 4. This means that $G$ of type {\bfseries pmm} has 4 Real structure:
one black-and-white of type 1,
two black-and-whites of type 2,
and finally one grey:
\[
(\mathrm{pmm} \leq \mathrm{p4m}), \ 
(\mathrm{pmm} \leq \mathrm{cmm}), \ 
(\mathrm{pmm} \leq \mathrm{pmm}), \
(\mathrm{pmm} \leq \mathrm{pmm}\times C_2).
\]
Simply listing the types of both $H$ and $G$ appears to be a convenient notation for the Real structures in the 2D case.
The only exception is that $(\mathrm{pm} \leq \mathrm{pm})$ may have two different meanings.
Therefore, we use  $(2\mathrm{pm} \leq \mathrm{pm})$ for one of them, where there are two subgroups conjugate by an automorphism.


In the row {\bfseries p2} and from the row {\bfseries pmm} down, 
the centre $Z(G)$ is trivial so that
all $\fR_{\theta}(G)$ are one point sets: these sizes are in the last column.
We also do not list individual members of $\Theta (G)$ -- we just list its size, equal to  ${\sN}$.
%we don't list individual 
%there we just list the cardinality of $\fR(G) = \cup_{\theta}\fR_{\theta}(G) $, automatically equal to $|\Theta (G)|$.
%
%$\fR(H) = \bigsqcup_{\theta \in \Out(H), \ \theta^2 =1} \fR_{\theta}(H)$ . 
%
In the type  {\bfseries p2}, the group $\Out (G)$ is infinite. Yet the conjugacy classes in $\Theta (G)$ are easy to describe -- 
here are their representatives:
\[ I_2, \; U_2, \; X_2, \; X_4, \; e_1, \; (X_2,e_1), \; (X_2,e_1+e_2). \]

In the four types with infinite centre, the right hand side of Table~\ref{table:wallpaper3}
is filled slightly differently, with more information. The third column lists representatives in $\Theta(G)$. 
Then we use Theorem~\ref{thm:EMLv1} and Proposition~\ref{prop:C2cohomology} to derive
\[ \fR_{\theta}(G) =
\HH^{2}(C_2, Z(G)_{\theta}) =
\dfrac{Z(G)^0_{\theta}}{(\theta + 1)Z(G)_{\theta}} \, , \]
subject to vanishing of the obstruction, which always vanishes in the 2D case.
For each individual $\theta \in \Theta (G)$, we list $Z(G)_{\theta}$ in the penultimate column, cf. $C_2$-modules in Proposition~\ref{prop:torsionfreeKS}.
The cardinality $|\fR_{\theta}(G)|$ is found by Formula~\eqref{eqn:cohomofspecialmodules} that, in our case, boils down to
  \begin{equation}
  \HH^{2}(\bZ_+) = \bZ/2\bZ, \quad
  \HH^{2}(\bZ_-) =
  \HH^{2}(\bZ_d) = 0 \, .
  \end{equation}
  Finally, we fill the cardinalities of $\Out_\theta (G)$-orbits on $\fR_{\theta}(G)$ into the last column.
  For example, consider the first row for {\bfseries p1}. Here $G=\bZ^2$. Since $I$ acts trivially, we need to add an element $t$ to $G$
  such that $tg=gt$ for all $g\in G$. The four elements of $\fR_{\theta}(G)$ correspond to four possible choices of $t^2$:
  \[ t^2 = 0 \ \mbox{ or } \ t^2 = e_1 \ \mbox{ or } \ t^2 = e_2 \ \mbox{ or } \ t^2 = e_1+e_2 \, . \]
  Clearly, the last three choices are all magnetic equivalent, that is, they comprise a single $\Out_\theta (G)$-orbit. We write $1+3$
  because $\fR_{\theta}(G)$ has  two orbits, with 1 and 3 elements.
